\section{Koko Eating Bananas}
\label{sec:bs-koko}

\problemstatement{There are $n$ piles of bananas, where pile $i$ has
$\textit{piles}[i]$ bananas. Koko can eat at most $k$ bananas from a single
pile in one hour (if a pile has fewer than $k$ bananas, she finishes that
pile in that hour and cannot start another pile in the same hour). Koko
wants to eat all the bananas within $h$ hours. Find the
\textbf{minimum} integer eating speed $k$ that allows her to finish within
$h$ hours.}

\begin{patternbox}
\emph{``Find the minimum speed/rate/capacity such that some task completes
within a time/resource limit''} is the archetypal \emph{binary search on
the answer} problem. The two keywords to look for together are
\textbf{``minimize the (speed/capacity/rate)''} and a \textbf{feasibility
constraint} (here, a deadline of $h$ hours) that becomes easier to satisfy
as the candidate answer increases. This combination -- minimize $x$ subject
to ``feasible($x$)'' becoming \texttt{true} once $x$ is large enough -- is
the single most common shape of binary-search-on-the-answer problems on
this sheet, and Koko Eating Bananas is its canonical introduction.
\end{patternbox}

\subsection*{Understanding the problem carefully}

For a fixed candidate speed $k$, the number of hours Koko needs to finish
pile $i$ is $\lceil \textit{piles}[i] / k \rceil$ (she needs a whole extra
hour for any remainder), so the total hours needed across all piles is
\[
  \textit{hoursNeeded}(k) = \sum_{i=1}^{n} \left\lceil
  \frac{\textit{piles}[i]}{k} \right\rceil.
\]
As $k$ increases, each term $\lceil \textit{piles}[i]/k \rceil$ can only
decrease or stay the same (eating faster never requires more hours for the
same pile), so $\textit{hoursNeeded}(k)$ is a non-increasing function of
$k$. Define the feasibility predicate
$P(k) = (\textit{hoursNeeded}(k) \le h)$. Since $\textit{hoursNeeded}$ is
non-increasing in $k$, $P$ is \texttt{false} for small $k$ and
\texttt{true} from some crossover point onward -- exactly the monotonic
structure binary search needs, with the answer being the smallest $k$
where $P$ becomes \texttt{true}.

\begin{ideabox}[Candidate Space and Predicate]
The candidate space is integer eating speeds $k \in [1, \max(\textit{piles})]$
(eating faster than the largest pile is never useful, since any pile is
finished in exactly one hour once $k$ reaches its size). The predicate is
$P(k) = (\textit{hoursNeeded}(k) \le h)$, and we want the smallest $k$
where $P$ holds.
\end{ideabox}

\subsection*{Why this reduces feasibility-checking to a simple scan}

The crucial design decision is to separate two concerns: \emph{deciding
which $k$ to try next} (binary search's job) and \emph{checking whether a
specific $k$ works} (a simple $O(n)$ linear scan over the piles, summing
up $\lceil \textit{piles}[i]/k \rceil$ for the fixed candidate $k$). This
separation is the heart of every binary-search-on-the-answer problem: we
never need to reason about the feasibility check's internal monotonicity
in $i$ -- only about the feasibility \emph{function as a whole}'s
monotonicity in the single candidate parameter $k$, which we established
above using a much simpler, purely algebraic argument (more bananas per
hour can never increase the hours needed).

\subsection*{Algorithm}

\begin{enumerate}
  \item Initialize $\textit{lo} \gets 1$,
        $\textit{hi} \gets \max(\textit{piles})$.
  \item While $\textit{lo} \le \textit{hi}$:
  \begin{enumerate}
    \item $\textit{mid} \gets \textit{lo}+(\textit{hi}-\textit{lo})/2$.
    \item Compute $\textit{hoursNeeded}(\textit{mid})$ by summing
          $\lceil \textit{piles}[i] / \textit{mid} \rceil$ over all piles.
    \item If $\textit{hoursNeeded}(\textit{mid}) \le h$: this speed works;
          record it and try slower: $\textit{hi} \gets \textit{mid}-1$.
    \item Else: too slow; speed up: $\textit{lo} \gets \textit{mid}+1$.
  \end{enumerate}
  \item Return the recorded answer (equivalently, $\textit{lo}$ once the
        loop ends).
\end{enumerate}

\begin{algorithm}[H]
\caption{Koko Eating Bananas}
\begin{algorithmic}[1]
\Procedure{MinEatingSpeed}{\textit{piles}, $h$}
  \State $\textit{lo} \gets 1$, $\textit{hi} \gets \max(\textit{piles})$
  \While{$\textit{lo} \le \textit{hi}$}
    \State $\textit{mid} \gets \textit{lo}+(\textit{hi}-\textit{lo})/2$
    \State $\textit{hours} \gets \sum_i \lceil \textit{piles}[i] / \textit{mid} \rceil$
    \If{$\textit{hours} \le h$}
      \State $\textit{hi} \gets \textit{mid} - 1$
    \Else
      \State $\textit{lo} \gets \textit{mid} + 1$
    \EndIf
  \EndWhile
  \State \Return \textit{lo}
\EndProcedure
\end{algorithmic}
\end{algorithm}

\subsection*{Dry run}

\dryrun{Take $\textit{piles} = [3, 6, 7, 11]$, $h = 8$.}

$\max(\textit{piles}) = 11$.

\begin{itemize}
  \item $\textit{lo}=1,\textit{hi}=11$: $\textit{mid}=6$. Hours:
        $\lceil3/6\rceil+\lceil6/6\rceil+\lceil7/6\rceil+\lceil11/6\rceil
        = 1+1+2+2=6 \le 8$. Feasible. $\textit{hi}=5$.
  \item $\textit{lo}=1,\textit{hi}=5$: $\textit{mid}=3$. Hours:
        $\lceil3/3\rceil+\lceil6/3\rceil+\lceil7/3\rceil+\lceil11/3\rceil
        =1+2+3+4=10 > 8$. Infeasible. $\textit{lo}=4$.
  \item $\textit{lo}=4,\textit{hi}=5$: $\textit{mid}=4$. Hours:
        $\lceil3/4\rceil+\lceil6/4\rceil+\lceil7/4\rceil+\lceil11/4\rceil
        =1+2+2+3=8 \le 8$. Feasible. $\textit{hi}=3$.
  \item $\textit{lo}=4 > \textit{hi}=3$: loop ends.
\end{itemize}

The answer is $\textit{lo} = 4$: the minimum eating speed is $4$ bananas
per hour.

\subsection*{C++ implementation}

\begin{lstlisting}
#include <bits/stdc++.h>
using namespace std;

long long hoursNeeded(vector<int>& piles, int k) {
    long long hours = 0;
    for (int p : piles) {
        hours += (p + k - 1) / k; // ceil(p / k)
    }
    return hours;
}

int minEatingSpeed(vector<int>& piles, int h) {
    int lo = 1, hi = *max_element(piles.begin(), piles.end());

    while (lo <= hi) {
        int mid = lo + (hi - lo) / 2;
        if (hoursNeeded(piles, mid) <= h) {
            hi = mid - 1;
        } else {
            lo = mid + 1;
        }
    }
    return lo;
}

int main() {
    vector<int> piles = {3, 6, 7, 11};
    cout << "Minimum eating speed: "
         << minEatingSpeed(piles, 8) << endl;
    return 0;
}
\end{lstlisting}

\begin{complexitybox}
\textbf{Time Complexity.} The candidate speed range $[1, \max(\textit{piles})]$
halves each iteration ($O(\log(\max(\textit{piles})))$ iterations), and
each feasibility check costs $O(n)$ to scan all piles, giving an overall
time complexity of
\[
  O(n \log(\max(\textit{piles}))).
\]
\textbf{Space Complexity.} $O(1)$ auxiliary space.
\end{complexitybox}

\begin{takeawaybox}
The general recipe for binary search on the answer: (1) identify the
quantity being minimized or maximized as the candidate; (2) write a
feasibility function that, for a fixed candidate, answers a yes/no
question with a simple (often linear) computation; (3) prove that
feasibility is monotonic in the candidate using simple algebraic reasoning
about the feasibility function's components; (4) binary search the
candidate range using that feasibility check as the predicate. Every
remaining ``binary search on the answer'' problem in this chapter follows
exactly this four-step recipe.
\end{takeawaybox}
